% from main.tex
%%%%%%%%%%  scan idx 96  |  folio 83  %%%%%%%%%%
équivaut à dire que \(F\) est pseudo-cohérent rel.\ à \(p\) et localement de tor-dimension finie rel.\ à \(p\) (i.e. rel.\ au faisceau d'anneaux \(p^{-1}(\mathcal O_S)\)).

Le théorème central de l'Exp.\ III est le \emph{théorème de finitude}, qui affirme (sous une forme un peu plus précise) que, si \(f:X\to Y\) est un morphisme \emph{propre} de \(S\)-schémas localement de type fini, le foncteur \(Rf_*\) transforme complexes pseudo-cohérents rel.\ à \(S\) en complexes pseudo-cohérents rel.\ à \(S\)~; ce théorème est en réalité une conjecture, mais est prouvé néanmoins dans les deux cas particuliers suivants : a)~\(S\) est localement noethérien; b)~\(f\) est projectif; il semble malheureusement que l'extension au cas général soit du même ordre de difficulté que le théorème analogue en Géométrie Analytique («~théorème~» de Grauert). Combinant le théorème de finitude avec une formule essentiellement triviale dite \emph{formule de projection}, on obtient des critères maniables pour la stabilité de la perfection relative par image directe. On retrouve, en corollaire, les «~théorèmes de continuité et de semi-continuité de Grauert~» (EGA III 7.6).

L'Exp.\ IV traduit en langage de «~groupes de Grothendieck~» les résultats des Exp.\ I à III. Sur un topos annelé \((X,\mathcal O_X)\), on note \(K^\bullet(X)\) (resp.\ \(K_\bullet(X)\)) le «~groupe de Grothendieck~» de la catégorie des complexes parfaits et de tor-dimension finie (resp.\ pseudo-cohérents et à cohomologie bornée). Le groupe \(K^\bullet(X)\) est un anneau, et \(K_\bullet(X)\) un module sur \(K^\bullet(X)\). \(K^\bullet\) est un foncteur contravariant (comme le suggère le point en haut), tandis que \(K_\bullet\) est un foncteur covariant pour les morphismes propres et pseudo-cohérents (de schémas ou d'espaces analytiques, et sous réserve que l'on dispose du théorème de finitude\dots). Il y a aussi des variances inhabituelles,
